% ch9_g 第9章 §9.7–§9.10 (书页 431–439 / p446–p454)
% 块边界判定：
%   起点——书页 431（p446）顶部为 9.9.1 节标题，起新小节，无前块溢出内容，
%     故本块不设 %JOIN%，亦无需 X-TAIL 区。
%   实际覆盖——9.9.1--9.9.5 节、9.10 节（含 9.10.1、9.10.2）及问题 9.10.1、9.10.2；
%     9.7、9.8 两节与 9.9 节引言在前块 ch9_f（书页 421--430）内。
%   终点——书页 439（p454）以问题 9.10.2 (c) 收束，为第 9 章末页（p455 为第 10 章
%     章首），自然收尾，无需 TAIL。
% 蓝本问题记录：
%   1) 式 (9.10.6) 右端蓝本与原书均印作 $U_{P_x}$，按上下文应为 $U_{P_y}$，已订正并加译注；
%   2) 式 (9.10.6) 引导句蓝本与原书均印作 $G_{P_y}T_{P_y}$，应为 $G_{P_y}P_y$，已按本意
%      译出并加译注；
%   3) 9.10.1 节蓝本 OCR 作 "IGG=\{c^{1(-)^{i}\theta\tau^{3}}\}"，按图片订正为
%      $\{\mathrm{e}^{\mathrm{i}(-)^{\pmb{i}}\theta\tau^{3}}\}$；
%   4) 9.10.2 节 "there will a gapless U(1) gauge boson" 原书漏 "be"，按本意译出；
%   5) 蓝本 9.9.4 节把 Z2Bx2(12)n 误排作 "Z2Bx2(12)n"（与图片一致，无碍）；
%      蓝本 9.9.3 节段首 "wefdw..." 等为文字层 OCR 噪声，未采信。

\subsection{投影对称群——量子相的一个普适性质}

\begin{keypoints}
\item 如果相应的平均场态是稳定的（即没有红外发散），PSG 就是真实量子相的一个普适性质。此时，PSG 描述该量子相中的量子序。
\end{keypoints}

在本节中，我们想证明 PSG 可以成为量子态的一种普适性质，其含义是它对微扰涨落是稳健的。因此，PSG 作为一种普适性质，可以用来刻画一个量子相。任何与 PSG 相联系（或受 PSG 保护）的物理性质也是相的普适性质，可以用于在实验中探测和测量量子序。特别地，PSG 能够保护无能隙规范激发和费米子激发（见 9.10 节）。因此，PSG 的稳定性也蕴含着无能隙规范激发和费米子激发的稳定性。

我们知道，平均场自旋液体态由 $U_{ij}=\langle\psi_{\pmb{i}}\psi_{\pmb{j}}^{\dagger}\rangle$ 刻画。如果在平均场态附近纳入微扰涨落，我们预期 $U_{ij}$ 会获得微扰修正 $\delta U_{ij}$。这里我们想论证：微扰涨落只能以这样的方式改变 $U_{ij}$，使得 $U_{ij}$ 与 $U_{ij}+\delta U_{ij}$ 具有相同的 PSG。

首先，我们想指出以下众所周知的事实。微扰涨落不能改变对称性和规范结构。例如，如果拟设 $U_{ij}$ 与哈密顿量具有某个对称性，那么由微扰涨落生成的 $\delta U_{ij}$ 将具有同样的对称性。类似地，微扰涨落不能生成这样的 $\delta U_{ij}$——比方说，把一个 $U(1)$ 规范结构破缺到 $Z_2$ 规范结构的 $\delta U_{ij}$。

由于规范结构（由 IGG 描述）和对称性群都是 PSG 的一部分，把上述观察推广为“除 PSG 中的 IGG 与对称性群之外，整个 PSG 都不能被微扰涨落改变”是合理的。

事实上，平均场哈密顿量和平均场基态在 PSG 中的变换下是不变的。因此，在围绕一个平均场态的微扰计算中，PSG 中的变换表现得就像普通的对称性变换一样。所以，微扰涨落只能生成在 PSG 的变换下不变的 $\delta U_{ij}$。

由于微扰涨落（按定义）不改变相，$U_{ij}$ 与 $U_{ij}+\delta U_{ij}$ 描述同一个相。换言之，我们可以把 $U_{ij}$ 分成若干类（称为普适类，universality classes），使得每一类中的 $U_{ij}$ 彼此通过微扰涨落相连。每一个普适类描述一个相。我们看到，如果上述论证成立，那么同一普适类中的拟设都享有同一个 PSG。换言之，普适类或者相是由 PSG（或量子序）分类的。

我们想指出，在上面的讨论中我们假定了微扰涨落没有红外发散。红外发散意味着微扰涨落是相关微扰。这种发散的修正可能引起相变，从而使上述论证失效。因此，上述论证和结果只适用于满足下列要求的稳定自旋液体：(i) 所有平均场涨落都没有红外发散；(ii) 平均场涨落在格点尺度上足够弱。

在下文中，我们将讨论几类平均场态的稳定性。要求 (ii) 可以通过大 $N$ 极限和/或（必要时）调整自旋哈密顿量中的短程自旋耦合来满足。这里我们将主要考虑要求 (i)。我们发现，至少在某些大 $N$ 极限下，许多（但并非全部）平均场态确实对应于在低能下稳定的真实量子自旋液体。

迄今为止所研究的所有自旋液体（每个元胞具有奇数个电子）可以分为四类。下面我们将逐一研究每一类。

\subsection{刚性自旋液体}

刚性自旋液体（rigid spin liquid）是自旋子和所有其他激发都完全有能隙的态。规范场的能隙可以由 Chern--Simons 项或 Anderson--Higgs 机制产生。有能隙的规范场只在自旋子之间诱导短程相互作用。由于低能下没有激发，刚性自旋液体是稳定的态。刚性自旋液体由拓扑序刻画，并且具有未禁闭的中性自旋-1/2 激发。刚性自旋液体的低能有效理论是拓扑场论。$Z_2$ 有能隙自旋液体和手征自旋液体是刚性自旋液体的例子。

\subsection{玻色自旋液体}

U1Cn00x 态 (9.8.9) 在式 (9.8.9) 中的 $a_{0}^{3}$ 足够大时，可以是一个 $U(1)$ 有能隙自旋液体。这样的态在平均场层次上具有有能隙的自旋子和无能隙的 $U(1)$ 规范玻色子。我们将称之为玻色自旋液体（Bose spin liquid）。无能隙 $U(1)$ 规范涨落的动力学由如下低能有效理论描述：
\begin{equation*}
\mathcal{L}=\frac{1}{8\pi g^{2}}\left(f_{\mu\nu}\right)^{2}
\end{equation*}
其中 $f_{\mu\nu}$ 是 $U(1)$ 规范场的场强。然而，在 $1+2$ 维中，并且在纳入瞬子效应之后，$U(1)$ 规范涨落将具有红外发散，这导致规范玻色子获得能隙以及自旋子被禁闭（Polyakov, 1977）。因此，平均场 $U(1)$ 有能隙态在 $1+2$ 维中是不稳定的。它们的 PSG 可能并不描述物理自旋液体中任何真实的量子序。\footnote{我们想\emph{提一句}，平均场 $U(1)$ 有能隙态在 $1+3$ 维中是\emph{稳定的}。}

\subsection{费米自旋液体}

费米自旋液体（Fermi spin liquid）由以下两条性质定义。其一，它们具有由自旋-1/2 费米子描述的无能隙激发；其二，这些无能隙激发之间只有短程相互作用。$Z_2$ 线性自旋液体、$Z_2$ 二次方自旋液体和 $Z_2$ 无能隙自旋液体是费米自旋液体的例子。

在 $Z_2$ 线性自旋液体中，自旋子具有无质量的狄拉克色散。这些具有短程相互作用的自旋子在连续极限下由如下有效理论描述（见问题 9.1.4）：
\begin{equation*}
S=\int\mathrm{d}t\,\mathrm{d}^{2}x\,\left[\bar{\psi}\partial_{\mu}\gamma^{\mu}\psi+\left(\bar{\psi}M\psi\right)^{2}\right]
\end{equation*}
在 $2+1$ 维中，$\psi$ 的标度量纲为 $[\psi]=1$，因此作用量是无量纲的。于是，相互作用项 $(\psi)^{4}$ 的标度量纲为 $[\psi^{4}]=4$，比 3 要大。我们看到，无质量狄拉克费米子之间的短程相互作用在 $1+2$ 维中是不相关的。因此，$Z_2$ 线性自旋液体是稳定的态。

现在让我们考虑式 (9.6.10) 中 $Z_2$ 二次方自旋液体 Z2Bx2(12)n 的稳定性。在 $Z_2$ 二次方自旋液体中，自旋子具有无能隙的二次方色散 $\omega\propto k^{2}$。在这种情形下，空间和时间具有不同的标度量纲，分别为 $[x^{-1}]\equiv1$ 与 $[t^{-1}]=2$。从无量纲的连续作用量
\begin{equation*}
S=\int\mathrm{d}t\,\mathrm{d}^{2}x\,\left[\psi^{\dagger}\mathrm{i}\partial_{t}\psi-c\psi^{\dagger}(\partial_{x})^{2}\psi\right]
\end{equation*}
我们看到，自旋子场 $\psi$ 的标度量纲为 $[\psi]=1$，而四费米子相互作用项的量纲为 $[\psi^{4}]=4$。于是，相互作用作用量 $S_{\mathrm{int}}=\int\mathrm{d}t\,\mathrm{d}^{2}x\,g(\psi^{\dagger}\psi)^{2}$ 中的耦合常数 $g$ 的标度量纲为 $[g]=0$，因为 $S_{\mathrm{int}}$ 的标度量纲总是零。因此，与 $Z_2$ 线性自旋液体不同，$Z_2$ 二次方态中无能隙自旋子之间的短程相互作用在 $1+2$ 维中是边缘的。相互作用的高阶效应会使耦合变得相关还是不相关，仍需进一步的研究才能确定。这将决定 $Z_2$ 二次方自旋液体在平均场层次之外的动力学稳定性。

在 $1+2$ 维中，$Z_2$ 无能隙自旋液体与费米液体一样稳定。如果我们假定费米液体在 $1+2$ 维中是稳定的，那么 $Z_2$ 无能隙自旋液体也是稳定的。

\subsection{代数自旋液体}

代数自旋液体（algebraic spin liquid）是具有无能隙激发的态，但其中没有任何一种无能隙激发是由自由的玻色型或费米型准粒子描述的。$U(1)$ 线性自旋液体是代数自旋液体的例子。它们的低能激发由与 $U(1)$ 规范场耦合的无质量狄拉克费米子描述。由于费米子--规范场耦合在微扰论的所有阶都是严格边缘的（Appelquist and Nash, 1990），无能隙激发一直低到零能都具有有限的相互作用。结果，低能下既没有自由的玻色型准粒子（如规范玻色子），也没有自由的费米型准粒子（如自旋子）。这使得关于这些态的稳定性的讨论困难得多。关于代数自旋液体的讨论，感兴趣的读者可参阅 Rantner and Wen (2002) 与 Wen (2002c)。

已经证明，在 $2+1$ 维中，$U(1)$ 线性态在自旋模型的某些平均场涨落微弱的大 $N$ 极限下是稳定的。$U(1)$ 线性态的 PSG 阻止微扰涨落生成破坏稳定性的抵消项（counter term），例如费米子质量项。这保证了 $U(1)$ 线性态的稳定性。因此，至少在大 $N$ 极限下，相应的代数自旋液体作为物理自旋系统的相而存在。

代数自旋液体的存在是一个非常惊人的现象。按照通常的看法，如果玻色子/费米子在低能下相互作用，那么相互作用将为这些低能激发打开一个能隙。这蕴含着：一个系统要么在低能下具有自由的玻色型/费米型激发，要么根本没有低能激发。代数自旋液体的存在表明，这种通常的看法是不正确的。它提出了一个重要问题：是什么保护着无能隙激发（特别是当它们在所有能标上都相互作用时）？无能隙激发的存在应当有一个“理由”或“原理”。在下一节中我们将证明，PSG 阻止规范玻色子和费米子获得质量项。\emph{所以，正是量子序及其相应的 PSG 保护着无能隙激发。}

\section{量子序与无能隙规范玻色子和费米子}

\begin{keypoints}
\item 量子序能够产生并保护无能隙规范玻色子与无能隙费米子，正如对称性破缺能够产生并保护无能隙 Nambu--Goldstone 玻色子一样。
\end{keypoints}

无能隙激发在自然界和凝聚态系统中都非常罕见。因此，如果我们看到一个无能隙激发，我们就会想问它为什么存在。无能隙激发的一种来源是自发对称性破缺，它给出 Nambu--Goldstone 玻色子（Nambu, 1960; Goldstone, 1961）。无能隙激发与自发对称性破缺之间的这种关系非常重要。正是由于这种关系，我们可以在不了解系统细节的情况下，仅从对称性出发得到一个复杂态的低能物理。这一思路使得 Landau 关于相与相变的对称性破缺理论（Ginzburg and Landau, 1950; Landau and Lifschitz, 1958）成为一个研究相的低能性质的极为强有力的理论。然而，自发对称性破缺并不是无能隙激发的唯一来源。量子序及其相应的 PSG 也能产生并保护无能隙激发。令人惊奇的是，量子序产生并保护的是无能隙规范玻色子和无能隙费米子。无能隙费米子甚至可以从纯玻色模型中产生。在本节中，我们想较为详细地讨论量子序（及其 PSG）与无能隙规范/费米子激发之间的关系（Wen, 2002a,c; Wen and Zee, 2002）。

\subsection{投影对称群与无能隙规范玻色子}

\begin{keypoints}
\item （在平均场层次上）无能隙规范玻色子的规范群就是 PSG 的 IGG。
\end{keypoints}

无能隙规范涨落与量子序之间的关系简单而直接。量子有序态中无能隙规范涨落的规范群，就是刻画该量子序的 PSG 中的 IGG。

为了看清量子序及其 PSG 如何产生并保护无能隙规范玻色子，作为一个例子，让我们假定一个量子有序态的 IGG $\mathcal{G}$ 包含一个 $U(1)$ 子群，它由下列常值规范变换构成：
\begin{equation*}
\{W_{\pmb{i}}=\mathrm{e}^{\mathrm{i}\theta\tau^{3}}\,|\,\theta\in[0,2\pi)\}\subset\mathcal{G}
\end{equation*}
接着我们考虑围绕平均场解 $\bar{u}_{ij}$ 的一类涨落，即 $u_{ij}=\bar{u}_{ij}\mathrm{e}^{\mathrm{i}a_{ij}^{3}\tau^{3}}$。由于 $\bar{u}_{ij}$ 在常值规范变换 $\mathrm{e}^{\mathrm{i}\theta\tau^{3}}$ 下不变，一个依赖空间的规范变换 $\mathrm{e}^{\mathrm{i}\theta_{\pmb{i}}\tau^{3}}$ 将把涨落 $a_{ij}^{3}$ 变为 $\tilde{a}_{ij}^{3}=a_{ij}^{3}+\theta_{\pmb{i}}-\theta_{\pmb{j}}$。这意味着 $a_{ij}^{3}$ 与 $\tilde{a}_{ij}^{3}$ 标记同一个物理态，而 $a_{ij}^{3}$ 对应于规范涨落。涨落的能量具有规范不变性，即 $E(\{a_{ij}^{3}\})=E(\{\tilde{a}_{ij}^{3}\})$。由于能量的这种规范不变性，我们看到规范场的质量项 $(a_{ij}^{3})^{2}$ 是不被允许的。因此，由 $a_{ij}^{3}$ 描述的 $U(1)$ 规范涨落将在低能下出现。

如果 $\mathcal{G}$ 的 $U(1)$ 子群由如下依赖空间的规范变换构成：
\begin{equation*}
\{W_{\pmb{i}}=\mathrm{e}^{\mathrm{i}\theta\,\pmb{n}_{\pmb{i}}\cdot\pmb{\tau}}\,|\,\theta\in[0,2\pi),\ |\pmb{n}_{\pmb{i}}|=1\}\subset\mathcal{G},
\end{equation*}
那么我们总可以用一个 $SU(2)$ 规范变换把每个格点上的 $\pmb{n}_{\pmb{i}}$ 转到 $z$ 方向，从而把问题约化为上面讨论过的情形。因此，无论 IGG 中的规范变换是否依赖空间，低能规范涨落的规范群总是由 $\mathcal{G}$ 给出。由于 IGG 的每一个 $U(1)$ 子群都对应于一个无能隙 $U(1)$ 规范玻色子，低能规范玻色子的规范群由 IGG 给出，即使 IGG 是非阿贝尔的。

我们想说明的是，有时低能规范涨落不仅出现在 $\pmb{k}=0$ 附近，还出现在其他一些 $\pmb{k}$ 点附近。在这种情形下，我们将有几个低能规范场，每个 $\pmb{k}$ 点对应一个。9.8 节中讨论的 $SU(2)$ 从属玻色子理论的某些拟设给出了这种现象的例子，它们在低能下具有 $SU(2)\times SU(2)$ 规范结构。我们看到，低能规范结构 $SU(2)\times SU(2)$ 甚至可以比高能规范结构 $SU(2)$ 还要大。即使对于低能规范涨落出现在不同 $\pmb{k}$ 点附近这种复杂情形，IGG 仍然正确地描述相应拟设的低能规范结构。如果 IGG 包含不依赖空间坐标的规范变换，那么这类变换对应于 $\pmb{k}=0$ 附近无能隙规范涨落的规范群。如果 IGG 包含依赖空间坐标的规范变换，那么这些变换对应于非零 $\pmb{k}$ 附近无能隙规范涨落的规范群。例如，若 $\mathrm{IGG}=\{\mathrm{e}^{\mathrm{i}(-)^{\pmb{i}}\theta\tau^{3}}\}$，则在 $\pmb{k}=(\pi,\pi)$ 附近就会出现一个无能隙 $U(1)$ 规范玻色子，由规范场 $(-)^{\pmb{i}}a^{3}(\pmb{x})$ 描述。这样，IGG 为我们提供了对所有低能规范涨落的统一处理，无论它们的晶体动量如何。

在本章中，我们在许多地方使用了 $Z_2$ 自旋液体、$U(1)$ 自旋液体、$SU(2)$ 自旋液体和 $SU(2)\times SU(2)$ 自旋液体这些术语。现在我们可以对这些低能 $Z_2$、$U(1)$、$SU(2)$ 和 $SU(2)\times SU(2)$ 规范群给出精确定义。这些低能规范群就是相应拟设的 IGG。它们与出现在 $SU(2)$、$U(1)$ 或 $Z_2$ 从属玻色子方法中的高能规范群毫无关系。我们还使用过一个平均场态的 $Z_2$ 规范结构、$U(1)$ 规范结构和 $SU(2)$ 规范结构这些术语。它们精确的数学含义同样是相应拟设的 IGG。当我们说一个 $U(1)$ 规范结构被破缺到一个 $Z_2$ 规范结构时，我们的意思是：拟设被改变了，使得它的 IGG 从 $U(1)$ 群变成了 $Z_2$ 群。

\subsection{投影对称群与无能隙费米子}

\begin{keypoints}
\item 有时，一个 PSG 保证无能隙费米子的存在。在这种情形下，PSG 的普适性意味着无能隙费米子的普适性（或稳定性）。
\end{keypoints}

为了演示 PSG 与无能隙费米子之间的直接联系，本节我们将研究一个特定的自旋液体，其量子序由 Z2Azz13 刻画（Wen and Zee, 2002）。Z2Azz13 自旋液体中的自旋子谱由自旋子跳跃哈密顿量
\begin{equation*}
H=\sum_{ij}\psi_{\pmb{i}}^{\dagger}u_{ij}\psi_{\pmb{j}}
\end{equation*}
给出。Z2Azz13 拟设的一个例子由式 (9.6.2) 给出。我们注意到，拟设 $u_{ij}$ 可以看作一个把费米子波函数 $\psi(\pmb{i})$ 映为 $\psi'(\pmb{i})=\sum_{\pmb{j}}u_{ij}\psi(\pmb{j})$ 的算符。这个算符记作 $\hat{H}$，将被称为哈密顿量。$\hat{H}$ 的本征值决定费米子谱。无能隙费米子对应于哈密顿量的零本征值。

Z2Azz13 的 PSG 由
\begin{align*}
G_{x}(\pmb{i}) &= \tau^{0}, & G_{y}(\pmb{i}) &= \tau^{0}\\
G_{P_{x}}(\pmb{i}) &= (-)^{\pmb{i}}\mathrm{i}\tau^{3} & G_{P_{y}}(\pmb{i}) &= (-)^{\pmb{i}}\mathrm{i}\tau^{3}\\
G_{P_{xy}}(\pmb{i}) &= \mathrm{i}\tau^{1} & G_{T}(\pmb{i}) &= \mathrm{i}\tau^{3}\\
G_{0}(\pmb{i}) &= -\tau^{0},\tag{9.10.1}
\end{align*}
生成。复合变换（如 $G_{P_{x}}P_{x}$）也可以看作作用在 $\psi_{\pmb{i}}$ 上的幺正算符。哈密顿量的投影对称性意味着，例如 $G_{P_{x}}P_{x}\hat{H}(G_{P_{x}}P_{x})^{\dagger}=\hat{H}$。

由于 $G_{x}=G_{y}=\tau^{0}$，Z2Azz13 拟设是平移不变的。在动量空间中，哈密顿量具有形式
\begin{equation*}
H(\pmb{k})=\epsilon^{\mu}(\pmb{k})\tau^{\mu}
\end{equation*}
拟设 $u_{ij}$ 在 $G_{T}T$ 下的不变性蕴含 $G_{T}(\pmb{i})u_{ij}G_{T}^{\dagger}(\pmb{j})=-u_{ij}$。用 $\pmb{k}$ 空间的语言表达，我们有
\begin{equation*}
U_{T}H(k_{x},k_{y})U_{T}^{\dagger}=-H(k_{x},k_{y}),\qquad U_{T}=\mathrm{i}\tau^{3}.\tag{9.10.2}
\end{equation*}
其中非平凡的 $G_{T}$ 给出非平凡的 $U_{T}$。式 (9.10.2) 蕴含 $\epsilon^{0,3}(\pmb{k})=0$，以及
\begin{equation*}
H(\pmb{k})=\epsilon^{1}(\pmb{k})\tau^{1}+\epsilon^{2}(\pmb{k})\tau^{2}\tag{9.10.3}
\end{equation*}
拟设在 $G_{P_{xy}}P_{xy}$ 下的不变性给出
\begin{equation*}
U_{P_{xy}}H(k_{x},k_{y})U_{P_{xy}}^{\dagger}=H(k_{y},k_{x}),\qquad U_{P_{xy}}=\mathrm{i}\tau^{1},\tag{9.10.4}
\end{equation*}
其中 $k_{x}$ 与 $k_{y}$ 的交换由 $P_{xy}$ 生成。拟设在 $G_{P_{x}}P_{x}$ 下的不变性导致
\begin{equation*}
U_{P_{x}}H(k_{x},k_{y})U_{P_{x}}^{\dagger}=H(-k_{x}+\pi,\,k_{y}+\pi),\qquad U_{P_{x}}=\mathrm{i}\tau^{3},\tag{9.10.5}
\end{equation*}
%FIG{9.16}: {$H(\pmb{k})$ 零点的两种图案。$\epsilon^{1}=0$ 线与 $\epsilon^{2}=0$ 线分别用实线和虚线表示。实心圆是缠绕数为 1 的零点，空心圆是缠绕数为 $-1$ 的零点。阴影区域是布里渊区的四分之一。当 $\epsilon^{2}$ 的零线移动而穿过 $\epsilon^{1}$ 的零线时，一个缠绕数为 1 的零点变成一个缠绕数为 $-1$ 的零点加上两个缠绕数为 1 的零点。}
而拟设在 $G_{P_{y}}P_{y}$ 下的不变性导致
\begin{equation*}
U_{P_{y}}H(k_{x},k_{y})U_{P_{y}}^{\dagger}=H(k_{x}+\pi,\,-k_{y}+\pi),\qquad U_{P_{y}}=\mathrm{i}\tau^{3}.\tag{9.10.6}
\end{equation*}
（译注：原书引导句误印作 $G_{P_{y}}T_{P_{y}}$，式 (9.10.6) 右端误印作 $U_{P_{x}}$，均按上下文订正为 $G_{P_{y}}P_{y}$ 与 $U_{P_{y}}$。）动量平移 $(\pi,\pi)$ 来自 $G_{P_{x}}$ 与 $G_{P_{y}}$ 中的 $(-)^{\pmb{i}}$ 因子。于是，非零的 $\epsilon^{1,2}(\pmb{k})$ 具有下列对称性：
\begin{align*}
\epsilon^{1}(k_{x},k_{y})&=-\epsilon^{1}(\pi-k_{x},\,\pi+k_{y})=-\epsilon^{1}(\pi+k_{x},\,\pi-k_{y})=\epsilon^{1}(k_{y},k_{x})\\
\epsilon^{2}(k_{x},k_{y})&=-\epsilon^{2}(\pi-k_{x},\,\pi+k_{y})=-\epsilon^{2}(\pi+k_{x},\,\pi-k_{y})=-\epsilon^{2}(k_{y},k_{x})\tag{9.10.7}
\end{align*}

事实上，式 (9.10.3) 与式 (9.10.7) 定义了 $\pmb{k}$ 空间中最一般的 Z2Azz13 拟设。

式 (9.10.7) 使我们能够由 $\epsilon^{1,2}(\pmb{k})$ 在布里渊区四分之一区域内的取值确定它们（见图 9.16）。在该四分之一布里渊区内，$\epsilon^{2}(\pmb{k})$ 在 $k_{x}$ 与 $k_{y}$ 互换下是反对称的。因此，当 $k_{x}=k_{y}$ 时 $\epsilon^{2}(\pmb{k})=0$。式 (9.10.7) 还蕴含 $\epsilon^{1}(\pmb{k})$ 在绕 $(0,\pi)$ 或 $(\pi,0)$ 的 $90^{\circ}$ 旋转下变号，即 $\epsilon^{1}(k_{x},k_{y}+\pi)=-\epsilon^{1}(k_{y},-k_{x}+\pi)$。因此，$\epsilon^{1}(\pmb{k})$ 必在 $(0,\pi)$ 与 $(\pi,0)$ 处为零。它在连接 $(0,\pi)$ 与 $(\pi,0)$ 的一条线上也必为零（见图 9.16）。结果，$\epsilon^{1}=0$ 线与 $\epsilon^{2}=0$ 线在四分之一布里渊区内至少相交一次。交点就是哈密顿量 $H(\pmb{k})$ 的零点。我们看到，无能隙自旋子至少出现在布里渊区中四个孤立的 $\pmb{k}$ 点上（见图 9.16）。拟设的 Z2Azz13 对称性直接导致无能隙自旋子。

$H(\pmb{k})$ 的零点的两类典型分布看上去可能就像图 9.16 中那样。我们看到，这些零点具有特殊的图案。为了理解图案中哪些特征是普适的，我们想研究当我们使拟设变形时零点的运动。在这样做之前，我们想指出，$H(\pmb{k})$ 的零点具有一种可以用缠绕数刻画的内部结构。当 $\pmb{k}$ 绕一个零点一周时，二维矢量 $(\epsilon^{1}(\pmb{k}),\epsilon^{2}(\pmb{k}))$ 就围绕 $(0,0)$ 画出一个回路。零点的缠绕数由该回路围绕 $(0,0)$ 缠绕的圈数给出。如果回路围绕 $(0,0)$ 逆时针缠绕，缠绕数为正；如果顺时针缠绕，则为负。典型的零点具有缠绕数 $\pm1$。比如说，一个缠绕数为 2 的零点，在我们微扰哈密顿量之后可以分裂成两个缠绕数为 1 的零点。

对于 Z2Azz13 拟设，$H(\pmb{k})$ 的零点由 $\epsilon^{1}(\pmb{k})$ 与 $\epsilon^{2}(\pmb{k})$ 的零线的交点给出。由于 $\epsilon^{0}(\pmb{k})=\epsilon^{3}(\pmb{k})=0$，所以当我们使拟设变形时，零点不能随意地出现或消失。零点只能成组地产生或湮没，并保持总缠绕数守恒（见图 9.16）。再结合拟设的对称性 (9.10.7)，我们发现下列性质对哈密顿量的任何不改变 PSG 的微小扰动都是稳健的：沿 $((0,0),(\pi,\pi))$ 连线的 $+1$ 与 $-1$ 零点的图案；以及 $(N_{+},N_{-})$，即三角形 $((0,0),(\pi,\pi),(0,\pi))$ \emph{内部} 的 $(+1,-1)$ 零点的数目（不包括边和角上的零点）。我们把这些性质定义为 Z2Azz13 自旋液体的零点图案（pattern of the zero's，POZ）。正如费米面拓扑的改变一样，POZ 的改变将导致基态能量出现奇异性，并标志一个相变。因此，POZ 也是刻画量子序的一个量子数。我们看到，自旋液体中的量子序并不能完全由 PSG 刻画。POZ 提供了一种额外的刻画。PSG 与 POZ 的结合为量子序提供了更完全的刻画。

\subsection*{问题 9.10.1}

考虑式 (9.6.1) 中 Z2A0013 态的一个拟设。你可以假定 $\gamma_{1,2}=\lambda=0$。

(a) 证明：改变 $a_{0}^{1}$ 将引起 POZ 的改变。证明 POZ 的改变会在平均场基态能量中引起奇异性。

(b) 求转变点处低能自旋子的谱。

\subsection*{问题 9.10.2}

(a) 求 Z2A003n PSG 的 $G_{x,y}$、$G_{P_{x},P_{y},P_{xy}}$ 与 $G_{T}$。

(b) 求 Z2A003n 态最一般的平均场哈密顿量 $H(\pmb{k})$。

(c) 证明：Z2A003n 态在 $\pmb{k}=(\pm\frac{\pi}{2},\pm\frac{\pi}{2})$ 处总是具有无能隙自旋子。

% ---------------- 新增术语（ch9_g） ----------------
% | rigid spin liquid | 刚性自旋液体 |
% | Bose spin liquid | 玻色自旋液体 |
% | Fermi spin liquid | 费米自旋液体 |
% | infra-red divergence | 红外发散 |
% | counter term | 抵消项 |
% | lattice scale | 格点尺度 |
% | spinon hopping Hamiltonian | 自旋子跳跃哈密顿量（hopping 系列词沿“跳跃”：跳跃问题/跳跃幅/跳跃算符） |
% | zero-line | 零线 |
% | pattern of the zero's (POZ) | 零点图案（POZ） |
% | winding number of a zero | 零点的缠绕数（缠绕数术语表已收） |
% | momentum shift | 动量平移 |
% | dynamical stability | 动力学稳定性 |
% | spatially-dependent (gauge transformation) | 依赖空间的（规范变换） |
% | constant gauge transformation | 常值规范变换 |
